
\medskip
\phantomsection\label{par:sga3-VIA-5.3}%
\noindent\textbf{5.3.---}
We now review several statements from Expos\'e~IV.

\medskip
\phantomsection\label{par:sga3-VIA-5.3.1}%
\noindent\textbf{5.3.1.---}
The statements of \SGARefLink{sga3:iv:proposition:5:2:7}{Expos\'e~IV, 5.2.7} and \SGARefLink{par:sga3-VIA-5.3.1}{5.3.1} translate as follows.
Let \(F\) and \(G\) be groups locally of finite type and flat over
\(A\), with \(F\) a closed normal subgroup of \(G\). The maps
\[
  H\longmapsto F\backslash H,
  \qquad
  H'\longmapsto H'\times_{(F\backslash G)}G
\]
define a bijection between the flat \(A\)-subgroups of \(G\)
containing \(F\) and the flat \(A\)-subgroups of \(F\backslash G\).
Under this bijection, the closed (respectively, normal) subgroups of
\(G\) containing \(F\) correspond to the closed (respectively,
normal) subgroups of \(F\backslash G\).%

\medskip
\phantomsection\label{par:sga3-VIA-5.3.2}%
\noindent\textbf{5.3.2.---}
\SGARefLink{sga3:iv:proposition:5:2:9}{Expos\'e~IV, Proposition~5.2.9}, gives the following result. Let
\(F,H,G\) be groups locally of finite type and flat over \(A\), and
suppose
\[
  F\subset H\subset G,
\]
with \(F\) closed in \(G\) and normal in \(H\). Then
\(F\backslash H\) acts freely on the left on \(F\backslash G\), the
quotient scheme
\[
  (F\backslash H)\backslash(F\backslash G)
\]
exists, and there is a canonical isomorphism of \(G\)-schemes
\[
  (F\backslash H)\backslash(F\backslash G)
    \xrightarrow{\ \sim\ }
  H\backslash G.
\]

\medskip
\phantomsection\label{par:sga3-VIA-5.3.3}%
\noindent\textbf{5.3.3.---}
Finally, \SGARefLink{sga3:iv:proposition:5:2:8}{Expos\'e~IV, 5.2.8}, gives the following assertion. Let
\(F,H,G\) be groups locally of finite type and flat over \(A\).
Suppose that \(F\) is contained in \(G\), closed and normal in \(G\),
that \(H\) is contained in \(G\), and that \(F\cap H\) is flat over
\(A\). Let
\[
  F\mathbin{\overset{\tau}{\times}}_A H
\]
be the \(A\)-group whose underlying scheme is \(F\times_AH\), with
multiplication defined by
\[
  \bigl((x,h),(y,h')\bigr)
    \longmapsto
  \bigl(xhyh^{-1},hh'\bigr).
\]
Likewise, let
\[
  u:H\cap F\longrightarrow
    F\mathbin{\overset{\tau}{\times}}_A H
\]
be the monomorphism \(x\mapsto(x^{-1},x)\), and let \(F\cdot H\) be
the quotient
\[
  (F\cap H)\backslash
    \bigl(F\mathbin{\overset{\tau}{\times}}_A H\bigr).
\]
Then there is a canonical isomorphism
\[
  F\backslash(F\cdot H)
    \xrightarrow{\ \sim\ }
  (F\cap H)\backslash H.
\]
